% !TEX root = ../../main.tex

\section{Permutaciones}

Una vez construido el grafo RAW, la extensión enumera ordenamientos topológicos del grafo. Cada ordenamiento topológico corresponde a una secuencia de statements que respeta todas las aristas de precedencia. Por lo tanto, cada uno representa una forma alternativa de escribir el bloque sin violar dependencias locales.

Para el ejemplo guía, el grafo tiene cinco ordenamientos válidos:

\begin{center}
\begin{tabular}{cl}
\hline
Candidato & Orden de statements \\
\hline
0 & \texttt{[1,2,3,4]} \\
1 & \texttt{[1,3,2,4]} \\
2 & \texttt{[1,3,4,2]} \\
3 & \texttt{[3,1,2,4]} \\
4 & \texttt{[3,1,4,2]} \\
\hline
\end{tabular}
\end{center}

El candidato 0 corresponde al orden original. Los demás candidatos mueven \texttt{s3} o \texttt{s1} sin violar las aristas RAW. Por ejemplo, \texttt{[1,3,2,4]} es válido porque \texttt{s1} sigue antes de \texttt{s2} y \texttt{s4}, y \texttt{s3} sigue antes de \texttt{s4}.

La enumeración de ordenamientos topológicos puede crecer rápidamente con la cantidad de statements independientes. Por ello, la extensión se presenta como experimental: permite estudiar el espacio de planes que se abre bajo conmutatividad, pero la explosión combinatoria debe considerarse una limitación y una línea de trabajo futuro.
